\documentclass[11pt]{article}

\usepackage{acl}
\usepackage{hyperref}
\usepackage{amsmath}
\usepackage{times}
\usepackage{latexsym}
\usepackage[T1]{fontenc}
\usepackage[utf8]{inputenc}
\usepackage{microtype}
\usepackage{inconsolata}
\usepackage{graphicx}
\usepackage{booktabs}

\title{Fact Omission Patterns Predict Media Bias: \\
A Discourse-Based Analysis of Cross-Ideological News Coverage}

\author{
Aditya Sai \\
IIIT Hyderabad \\
\url{heyadityawork@gmail.com}
\And
Aayush Batra \\
IIIT Hyderabad \\
\url{aayushbat06@gmail.com}
\And
Yash More \\
IIIT Hyderabad \\
\url{yashsmore7499@gmail.com}
}

\begin{document}
\maketitle

\begin{abstract}
We investigate whether the structural prominence of facts within Rhetorical Structure Theory (RST) discourse trees can predict the ideological lean of news articles. Using triplets of ideologically diverse articles (left, center, right) covering the same events, we align shared facts across articles via semantic clustering and compute Document Framing Index (DFI) features based on RST depth and nuclearity. Contrary to our initial hypothesis, we find that \textbf{RST structural placement alone provides no predictive signal}---models trained on structural features without coverage information achieve only 51\% accuracy (near chance). Instead, the dominant signal is \textbf{fact omission}: which facts a source chooses to cover or omit relative to a centrist baseline. A Random Forest classifier using bipartite EDU matching achieves \textbf{89.77\%} accuracy in distinguishing left-vs-center from right-vs-center samples. Analysis reveals that right-leaning sources systematically omit more facts that centrist sources cover, particularly in early (salient) positions. Our work provides a novel, interpretable lens for media bias analysis: bias manifests primarily through selective fact coverage rather than rhetorical positioning of covered facts.
\end{abstract}

\section{Introduction}

Media bias is a widely studied phenomenon in journalism and computational social science. While bias is often associated with explicit opinionated language, many instances of bias emerge through more subtle mechanisms of framing and emphasis. News outlets may describe the same event while selectively highlighting certain facts, presenting supporting information as central content, or relegating competing perspectives to background context. Such structural differences influence how readers interpret events even when the underlying factual content is similar.

Most existing approaches to automated media bias detection rely on lexical or stylistic signals. Recent work using transformer-based models has achieved strong performance in classifying ideological leaning \citep{Menzner_2024}. However, these models operate largely as black boxes, offering limited insight into \emph{why} an article appears biased.

Discourse structure provides a promising lens for studying these phenomena. Rhetorical Structure Theory \citep{nguyen2021rstparsingscratch} models documents as hierarchical trees of elementary discourse units (EDUs), where relations between units specify how pieces of information support, elaborate on, or contrast with one another. Importantly, RST distinguishes between \emph{nucleus} segments, which represent central information, and \emph{satellite} segments, which provide supporting material. This hierarchical representation naturally encodes the prominence of information within a document.

In this work, we investigate whether this structural prominence can predict ideological bias. We construct triplets of articles covering the same news event from left, center, and right sources, align shared facts across articles using semantic clustering, and compare how these facts are positioned within each article's RST tree.

\paragraph{Research Question} Can the structural prominence of a fact in an RST tree---measured using depth and nuclearity in the discourse hierarchy---predict the ideological lean of a news outlet?

\paragraph{Key Findings} Our experiments reveal a surprising answer: \textbf{RST structural placement alone cannot predict bias}. When we control for fact coverage by restricting analysis to facts covered by all three ideological sources, classification accuracy drops to 51\%---indistinguishable from chance. Instead, the dominant predictive signal is \textbf{which facts are omitted}: different bias orientations systematically cover or omit different facts, and this coverage pattern alone achieves strong classification performance. Our best approach---bipartite EDU matching that explicitly encodes coverage asymmetry---achieves 89.77\% accuracy with Random Forest.

\paragraph{Contributions}
\begin{enumerate}
    \item A novel dataset of 1,730 fact-aligned triplets with RST annotations, enabling discourse-level bias analysis.
    \item Empirical evidence that RST structural prominence does not predict bias when coverage is controlled.
    \item Discovery that fact omission patterns---particularly asymmetric omission rates between left and right sources---are the dominant signal for bias classification.
    \item Analysis of what patterns classifiers learn from unaligned DFI vectors, revealing distributional regularities that generalize across events.
\end{enumerate}


\section{Background and Related Work}

\subsection{Lexical and Informational Bias}

Media bias manifests in multiple forms and has been systematically categorized in recent literature. \citet{spinde2024mediabiastaxonomysystematic} review thousands of studies on automated bias detection and introduce the \emph{Media Bias Taxonomy}, which organizes different types of bias and surveys computational approaches used to detect them.

Within this landscape, a key distinction is between \emph{lexical bias} and \emph{informational bias}. \citet{fan2019plainsightmediabias} introduce the BASIL dataset and show that informational bias---factual content that nonetheless shapes reader perception through selective presentation---appears frequently in news articles and poses challenges for automatic detection. Building on this, \citet{berg2020contextinformationalbiasdetection} demonstrate that detecting informational bias benefits from incorporating broader context, highlighting the inherently context-dependent nature of this bias type.

\subsection{Transformer-Based Bias Classification}

Recent work applies transformer-based models such as BERT and RoBERTa to detect political bias in news text. \citet{banik2025bridginghumanmodelperspectives} compare human annotations with predictions from multiple models to analyze how closely automated systems align with human judgments of political slant. \citet{ghosh2025explainingnewsbiasdetection} investigate the internal decision mechanisms of transformer-based bias detectors using SHAP explanations, showing that different architectures rely on evaluative language cues differently. Together, these studies highlight both the effectiveness of transformer models for bias detection and the importance of interpretability.

\subsection{Discourse-Based Bias Analysis}

A growing body of work explores the connection between discourse structure and perceived bias. \citet{asher2018biassemanticdiscourseinterpretation} argue that interpretive bias emerges from the interaction between discourse structure and conversational dynamics, showing how the role a sentence plays in the discourse can influence how strongly its content is interpreted as biased.

However, these methods operate at the level of individual sentences and rely on surface-level positional cues. Our approach uses full RST trees \citep{nguyen2021rstparsingscratch}, capturing hierarchical relationships of nuclearity and depth to quantify how strongly a fact is emphasized or downplayed across entire documents.

\subsection{Rhetorical Structure Theory}

RST represents a document as a constituency tree over EDUs. Internal nodes encode a rhetorical relation (e.g., \textsc{Elaboration}, \textsc{Evidence}, \textsc{Contrast}) and a nuclearity assignment: one child is the \emph{nucleus} (the more essential content) and the other the \emph{satellite} (supporting material). Nucleus EDUs closer to the root are structurally prominent; satellite EDUs deep in the tree are structurally marginal.


\section{Dataset}

\subsection{Source Corpus}

We use 37,554 news articles labeled by political lean (left, center, right) from the Article-Bias-Prediction dataset\footnote{\url{https://github.com/ramybaly/Article-Bias-Prediction}}, which contains articles scraped from AllSides.com. Each article includes fields for article ID, source outlet, title, content, and a bias label.

\subsection{Triplet Construction}

The dataset does not natively group articles by news event. To identify \emph{triplets}---groups of one left, one center, and one right article covering the same event---we applied an unsupervised nearest-neighbor retrieval procedure.

Each article was embedded using the \texttt{all-MiniLM-L6-v2} sentence transformer (first 1,000 characters only). For each article, the 10 nearest neighbors were retrieved using cosine similarity. All combinations of three mutually similar articles (similarity $> 0.8$) whose bias labels form $\{\text{left}, \text{center}, \text{right}\}$ were stored as triplets.

This procedure yielded 6,714 raw triplets. After removing duplicates (from anchor-based enumeration) and applying a participation cap of 5 triplets per article to reduce over-representation, we retained \textbf{2,631 unique triplets}.

\subsection{RST Parsing and Valid Subset}

We use the \texttt{isanlp\_rst\_v3} neural discourse parser to generate RST trees for each article, with a 5-second timeout per document. RST trees were successfully generated for 3,005 articles.

Excluding triplets containing articles without RST parses yields \textbf{1,748 valid triplets} for downstream analysis. These triplets span 2,617 unique articles.

\subsection{Leakage-Safe Data Split}

We split the 1,748 valid triplets into train/validation/test (80/10/10) under a hard constraint: no article may appear across different splits. This prevents information leakage through shared articles.

The split yields:
\begin{itemize}
    \item Train: 1,398 triplets (2,766 binary samples)
    \item Validation: 174 triplets (342 samples)
    \item Test: 176 triplets (352 samples)
\end{itemize}


\section{Methodology}

\subsection{Cross-Document Fact Alignment}

To compare rhetorical framing across articles, we identify when different documents refer to the same underlying fact using a two-stage semantic clustering pipeline.

\paragraph{Stage 1: SBERT Clustering}
Each EDU is embedded using SBERT (\texttt{all-mpnet-base-v2}). All EDUs from a triplet are pooled and clustered using agglomerative clustering with cosine distance and a threshold of 0.35. Clusters containing EDUs from multiple articles are candidate fact alignments.

\paragraph{Stage 2: Cross-Encoder Refinement}
To reduce false positives, we apply cross-encoder reranking (\texttt{cross-encoder/ms-marco-MiniLM-L-6-v2}). Only EDU pairs scoring above 0.7 are accepted. Clusters must contain at least one EDU from the center article (our reference baseline).

\paragraph{EDU Filtering}
Before clustering, we filter non-factual EDUs: those with $<3$ tokens, punctuation-only fragments, URLs, boilerplate phrases (``Story highlights''), and short attribution shells (``he said'').

This pipeline produces \textbf{16,170 validated fact clusters} across 1,730 triplets.

\subsection{Document Framing Index (DFI)}

For each fact cluster $c$, we compute a prominence score for each bias side based on RST tree position. If a side has multiple EDUs in a cluster, we take the maximum prominence.

\paragraph{Original Prominence Formula}
\begin{equation}
W(d, s) = \alpha^{d+1} \cdot \gamma^s
\end{equation}
where $d$ is RST tree depth, $s$ is the count of satellite edges to root, and $\alpha, \gamma \in (0,1)$ are decay parameters.

\paragraph{Log-Depth Formula (Best Performing)}
\begin{equation}
W(d) = \frac{1}{1 + \ln(1+d)}
\end{equation}

The DFI delta for each cluster encodes the prominence difference relative to center:
\begin{equation}
\Delta^{(c)}_{\text{side}} = W_{\text{side}}(c) - W_{\text{center}}(c)
\end{equation}

Positive values indicate the side \emph{amplifies} that fact; negative values indicate \emph{suppression}; zero for omitted facts (when a side has no EDU in that cluster).

\subsection{Feature Representations}

We evaluate three feature encodings:

\paragraph{Coverage Features}
Binary presence/absence deltas: $h_{\text{side}} - h_{\text{center}} \in \{-1, 0, +1\}$, where $h=1$ if the side has at least one EDU in the cluster. This encodes \emph{which facts are omitted}.

\paragraph{Structural Features}
Log-depth prominence deltas: continuous values encoding \emph{how prominently} covered facts are positioned in the RST tree.

\paragraph{Combined Features}
Both coverage and structural features interleaved per cluster (2 features per cluster, 138 total dimensions).

All features are zero-padded to a fixed length (69 dimensions for single features, 138 for combined).

\subsection{Classification Task}

Each triplet contributes two binary samples:
\begin{itemize}
    \item Left-vs-center (label 0): Is this sample from the left article?
    \item Right-vs-center (label 1): Is this sample from the right article?
\end{itemize}

This formulation ensures exact class balance and frames the task as: given a DFI vector representing how a non-centrist source differs from center, predict whether it is left or right.


\section{Experiments}

\subsection{Experiment 1: Structural Ablation}

We first test whether RST structural parameters contribute predictive signal.

\begin{table}[h]
\centering
\small
\begin{tabular}{lcc}
\toprule
\textbf{Configuration} & \textbf{Val Acc} & \textbf{Test Acc} \\
\midrule
Baseline ($\alpha$=0.8, $\gamma$=0.5) & 69.88\% & 67.05\% \\
Without satellite ($\gamma$=1) & 68.13\% & 70.17\% \\
Without depth ($\alpha$=1) & 71.35\% & 65.91\% \\
Without both ($\alpha$=1, $\gamma$=1) & \textbf{78.65\%} & \textbf{75.28\%} \\
\bottomrule
\end{tabular}
\caption{Structural parameter ablation with SVM. Removing \emph{both} structural parameters yields the best performance, suggesting structural weighting adds noise.}
\label{tab:structural_ablation}
\end{table}

\textbf{Finding:} Removing both depth and satellite weighting \emph{improves} accuracy from 67\% to 75\%. This indicates that structural parameters introduce noise rather than signal.

\subsection{Experiment 2: Coverage-Only Baseline}

The ``without both'' condition is mathematically equivalent to binary coverage features. We verify this explicitly:

\begin{table}[h]
\centering
\small
\begin{tabular}{lcc}
\toprule
\textbf{Model} & \textbf{Test Acc} & \textbf{Test F1} \\
\midrule
Structural baseline & 67.05\% & 67.02\% \\
Coverage-only & \textbf{75.28\%} & \textbf{75.04\%} \\
\bottomrule
\end{tabular}
\caption{Coverage features outperform structural features by 8 percentage points.}
\label{tab:coverage_baseline}
\end{table}

\textbf{Finding:} Coverage-only features achieve 75.28\%---identical to ``without both''---confirming that the model learns from omission patterns, not RST structure.

\subsection{Experiment 3: Coverage-Controlled Analysis}

To isolate structural signal, we restrict to clusters where all three sides are present (``size-3'' clusters), removing coverage variation.

\begin{table}[h]
\centering
\small
\begin{tabular}{lcc}
\toprule
\textbf{Configuration} & \textbf{Test Acc} & \textbf{Test F1} \\
\midrule
Structural (all clusters) & 76.42\% & 76.41\% \\
Structural (size-3 only) & 51.42\% & 48.70\% \\
\bottomrule
\end{tabular}
\caption{When coverage is controlled, structural features drop to chance level.}
\label{tab:size3_ablation}
\end{table}

\textbf{Finding:} With coverage removed, accuracy drops to 51\%---near chance. This definitively shows that \textbf{RST structural placement alone cannot predict bias}.

\subsection{Experiment 4: Alternative Prominence Formulas}

We test whether alternative structural formulas can improve upon the original:

\begin{table}[h]
\centering
\small
\begin{tabular}{lcc}
\toprule
\textbf{Formula} & \textbf{Test Acc} & \textbf{vs Coverage} \\
\midrule
Original $\alpha^{d+1}\gamma^s$ & 67.05\% & $-8.23\%$ \\
Exponential $e^{-0.3d}$ & 67.05\% & $-8.23\%$ \\
Inverse satellite $1/(1+0.5s)$ & 75.00\% & $-0.28\%$ \\
Log-depth $1/(1+\ln(1+d))$ & \textbf{76.42\%} & $+1.14\%$ \\
Coverage-only & 75.28\% & --- \\
\bottomrule
\end{tabular}
\caption{Alternative prominence formulas. Log-depth slightly outperforms coverage-only.}
\label{tab:formulas}
\end{table}

\textbf{Finding:} Log-depth achieves 76.42\%, marginally above coverage-only (+1.14\%). This is the only structural variant that outperforms coverage.

\subsection{Experiment 5: Combined Features}

We test whether combining coverage and structural features yields improvement:

\begin{table}[h]
\centering
\small
\begin{tabular}{lccc}
\toprule
\textbf{Features} & \textbf{Dim} & \textbf{Test Acc} & \textbf{Test F1} \\
\midrule
Coverage-only & 69 & 75.28\% & 75.04\% \\
Structural (log-depth) & 69 & 76.42\% & 76.41\% \\
Combined & 138 & \textbf{77.27\%} & \textbf{77.27\%} \\
\bottomrule
\end{tabular}
\caption{Combined features achieve the best SVM performance.}
\label{tab:combined}
\end{table}

\textbf{Finding:} Combined features achieve 77.27\%, suggesting coverage and structural signals provide modest complementary information.

\subsection{Experiment 6: Model Comparison}

We compare multiple model families:

\begin{table}[h]
\centering
\small
\begin{tabular}{lccc}
\toprule
\textbf{Model} & \textbf{Coverage} & \textbf{Structural} & \textbf{Combined} \\
\midrule
SVM (RBF) & 76.70\% & 77.56\% & 77.84\% \\
Logistic (L2) & 74.43\% & 73.58\% & 75.28\% \\
MLP & 79.26\% & 78.12\% & 76.70\% \\
\textbf{Random Forest} & 79.26\% & 86.08\% & \textbf{86.36\%} \\
\bottomrule
\end{tabular}
\caption{Model comparison. Random Forest dramatically outperforms other models.}
\label{tab:models}
\end{table}

\textbf{Finding:} Random Forest achieves 86.36\% on combined features---8.5 points higher than SVM.

\subsection{Experiment 7: Hyperparameter Tuning}

Grid search on Random Forest with 5-fold cross-validation:

\begin{table}[h]
\centering
\small
\begin{tabular}{lcc}
\toprule
\textbf{Model} & \textbf{CV Acc} & \textbf{Test Acc} \\
\midrule
Random Forest (default) & 84.12\% & 86.36\% \\
Random Forest (tuned) & 85.20\% & \textbf{87.78\%} \\
\bottomrule
\end{tabular}
\caption{Hyperparameter tuning results. Best: depth=None, n\_estimators=50, min\_samples\_split=5.}
\label{tab:tuning}
\end{table}

\textbf{Final Result:} Tuned Random Forest achieves \textbf{87.78\% test accuracy}.

\subsection{Experiment 8: Cross-Validation and Confidence Intervals}

To assess stability, we run 5-fold cross-validation:

\begin{table}[h]
\centering
\small
\begin{tabular}{lccc}
\toprule
\textbf{Features} & \textbf{Mean Acc} & \textbf{Std} & \textbf{95\% CI} \\
\midrule
Coverage & 77.34\% & 0.85\% & [75.4\%, 81.7\%] \\
Structural & 77.02\% & 0.90\% & [73.3\%, 79.3\%] \\
Combined & 76.65\% & 1.35\% & [72.4\%, 78.8\%] \\
\bottomrule
\end{tabular}
\caption{5-fold cross-validation results with SVM.}
\label{tab:cv}
\end{table}

\textbf{Finding:} Results are stable across folds (std $<1.5\%$). The overlapping confidence intervals suggest differences between feature sets may not be statistically significant with SVM.

\subsection{Experiment 9: Aggregate Feature Vectors}

A conceptual concern with the padded DFI approach is that different triplets have different numbers of clusters, meaning the $i$-th dimension corresponds to different facts across samples. Padding and concatenating these vectors mixes unrelated dimensions.

We test an alternative: extract fixed-length \textbf{aggregate statistical features} from each DFI vector, making dimensions semantically consistent across samples.

\paragraph{Aggregate Features (9 per side)}
\begin{itemize}
    \item Mean, std, min, max of DFI deltas
    \item Skewness (asymmetry of distribution)
    \item Proportion positive deltas (side $>$ center)
    \item Proportion zero deltas (equal coverage)
    \item Proportion negative deltas (center $>$ side)
    \item Log cluster count: $\ln(1 + K)$
\end{itemize}

\begin{table}[h]
\centering
\small
\begin{tabular}{lccc}
\toprule
\textbf{Features} & \textbf{Model} & \textbf{Dim} & \textbf{Test Acc} \\
\midrule
Aggregate coverage & SVM & 9 & 60.51\% \\
Aggregate coverage & RF & 9 & 54.83\% \\
Aggregate combined & SVM & 18 & 61.36\% \\
Aggregate combined & RF & 18 & 55.11\% \\
Aggregate (norm) & RF & 9 & 65.06\% \\
\midrule
Padded DFI & RF & 138 & \textbf{87.78\%} \\
\bottomrule
\end{tabular}
\caption{Aggregate features vs.\ padded DFI. Aggregate features fail dramatically.}
\label{tab:aggregate}
\end{table}

\textbf{Finding:} Aggregate features achieve at most 65\% accuracy---far below padded DFI (87.78\%). The high validation accuracy ($>$97\%) with poor test performance indicates severe overfitting. This reveals that \textbf{positional information is critical}: the model learns that position 0 (first cluster) behaves differently from position 5, even though these correspond to different facts across triplets. Aggregate statistics destroy this position-specific signal.

\subsection{Experiment 10: Alternative DFI Construction Methods}

We test three alternative approaches to constructing DFI feature vectors from clustered facts.

\paragraph{Alternative 1: Cumulative Prominence}
Replace $W_{\text{doc}} = \max(\text{scores})$ with $W_{\text{doc}} = \sum(\text{scores})$ or $W_{\text{doc}} = \text{mean}(\text{scores})$. This captures both structural prominence \emph{and} repetition frequency---if a fact is mentioned 10 times with low prominence, its cumulative score still rises.

\paragraph{Alternative 2: Distributional DFI}
Expand each cluster from a 1D scalar to a 3D or 4D vector:
\begin{itemize}
    \item $\max(W)$: Peak prominence (headline placement)
    \item $\sum(W)$: Total structural emphasis
    \item $\text{count}$: Pure repetition frequency
    \item (Optional) Coverage: Binary presence
\end{itemize}
This lets the classifier learn which aspect matters most for different bias patterns.

\paragraph{Alternative 3: Bipartite Decomposition}
Instead of comparing ``best of 5 vs.\ best of 1'' when one side has more mentions, perform 1-to-1-to-1 matching of EDUs across sides using structural similarity. Leftover EDUs become explicit omission signals (matched against 0). This directly encodes coverage asymmetry.

\begin{table}[h]
\centering
\small
\begin{tabular}{lccc}
\toprule
\textbf{Method} & \textbf{Variant} & \textbf{Dim} & \textbf{Test Acc} \\
\midrule
\multicolumn{4}{l}{\textit{Baseline (current approach)}} \\
max(prominence) & structural & 69 & 87.50\% \\
coverage-only & binary & 69 & 80.40\% \\
combined & cov + struct & 138 & 87.22\% \\
\midrule
\multicolumn{4}{l}{\textit{Alternative 1: Cumulative Prominence}} \\
sum(prominence) & unnormalized & 69 & 82.67\% \\
avg(prominence) & normalized & 69 & 88.07\% \\
coverage + sum & combined & 138 & 87.78\% \\
\midrule
\multicolumn{4}{l}{\textit{Alternative 2: Distributional DFI}} \\
[max, sum, count] & 3D & 207 & 86.08\% \\
[cov, max, sum, count] & 4D & 276 & 86.36\% \\
[max, avg, log\_count] & normalized & 207 & 86.65\% \\
\midrule
\multicolumn{4}{l}{\textit{Alternative 3: Bipartite Decomposition}} \\
1-to-1 matching & prominence & 157 & 88.07\% \\
1-to-1 matching & coverage & 157 & \textbf{89.77\%} \\
1-to-1 matching & combined & 314 & 88.07\% \\
\bottomrule
\end{tabular}
\caption{DFI construction alternatives with Random Forest. Bipartite decomposition with coverage features achieves the best performance.}
\label{tab:alternatives}
\end{table}

\textbf{Key Findings:}

\begin{enumerate}
    \item \textbf{Alternative 3 (Bipartite) achieves best performance}: The coverage-based bipartite decomposition reaches \textbf{89.77\%}---a +2.0\% improvement over the baseline. This approach explicitly decomposes ``5 mentions vs.\ 1 mention'' into 1 matched pair plus 4 explicit omissions.
    
    \item \textbf{Alternative 1 shows modest improvement}: Using $\text{avg}(\text{prominence})$ instead of $\text{max}(\text{prominence})$ achieves 88.07\%, capturing the repetition signal currently lost by taking only the maximum.
    
    \item \textbf{Alternative 2 underperforms}: Despite richer representation (3--4 features per cluster), distributional DFI performs slightly below baseline. The increased dimensionality may dilute the key signal or cause mild overfitting.
\end{enumerate}

The success of bipartite decomposition confirms our core finding: \textbf{coverage asymmetry is the dominant signal}. By explicitly encoding each unmatched EDU as an omission event, this approach maximizes the information available about selective fact coverage.


\section{Analysis: What Does the Model Learn?}

A key question arises: each triplet has different fact clusters, so the $i$-th dimension in one DFI vector corresponds to a completely different fact than in another vector. Yet Random Forest achieves up to 89.77\% accuracy. What patterns does it learn?

\subsection{Pattern 1: Asymmetric Omission Rates}

The most robust pattern is that \textbf{right-leaning sources omit more facts that center covers}:

\begin{table}[h]
\centering
\small
\begin{tabular}{lcc}
\toprule
\textbf{Sample Type} & \textbf{Mean Omissions} & \textbf{Std} \\
\midrule
Left omits, center covers & 3.52 & 4.68 \\
Right omits, center covers & 3.97 & 4.31 \\
\bottomrule
\end{tabular}
\caption{Right sources omit $\approx$13\% more facts than left sources.}
\label{tab:omission}
\end{table}

\subsection{Pattern 2: Position-Dependent Distributions}

Early positions show stronger systematic differences:

\begin{table}[h]
\centering
\small
\begin{tabular}{cccc}
\toprule
\textbf{Position} & \textbf{Left Mean} & \textbf{Right Mean} & \textbf{Diff} \\
\midrule
0 & $-0.069$ & $-0.566$ & $-0.497$ \\
1 & $-0.159$ & $-0.504$ & $-0.345$ \\
2 & $-0.222$ & $-0.452$ & $-0.230$ \\
3 & $-0.244$ & $-0.422$ & $-0.177$ \\
\bottomrule
\end{tabular}
\caption{Coverage delta means by position. Early positions show larger differences.}
\label{tab:positions}
\end{table}

\subsection{Why Random Forest Excels}

\begin{enumerate}
    \item \textbf{Automatic feature selection}: Trees focus on early positions where signal is strongest. Position 0 accounts for 39.4\% of total importance.
    \item \textbf{Threshold-based splits}: Trees naturally handle discrete coverage deltas $\{-1, 0, +1\}$.
    \item \textbf{Robustness to sparse features}: Most DFI values are zero (padding). Trees efficiently ignore uninformative zeros.
\end{enumerate}

\subsection{What the Model Does NOT Learn}

\begin{itemize}
    \item \textbf{Not semantic fact content}: The model has no knowledge of what facts are about.
    \item \textbf{Not cross-triplet fact alignment}: Cluster 3 in triplet A may be completely different from cluster 3 in triplet B.
    \item \textbf{Not RST structural placement}: Size-3 ablation shows structural signals alone achieve only 51\%.
\end{itemize}

\subsection{The Learned Pattern}

The model learns a \textbf{distributional fingerprint}:

\begin{quote}
\textit{Right-leaning sources systematically omit more facts that centrist sources cover, and this omission pattern is particularly pronounced in the first few (most salient) facts of each event.}
\end{quote}


\section{Discussion}

\subsection{Implications for Bias Research}

Our findings suggest a reframing of discourse-based bias analysis. Rather than ``RST structure predicts bias,'' the more accurate statement is: \textbf{fact omission patterns predict bias}. This has several implications:

\begin{enumerate}
    \item \textbf{Coverage matters more than positioning}: How facts are rhetorically positioned is less predictive than whether they are covered at all.
    
    \item \textbf{Interpretability}: Omission-based features are more interpretable than RST structural features. We can directly identify which facts are omitted by which sources.
    
    \item \textbf{Model selection matters}: Random Forest's 10-point improvement over SVM suggests that model choice can dominate feature engineering in this task.
\end{enumerate}

\subsection{Why Does RST Structure Not Predict Bias?}

Several factors may explain this:

\begin{enumerate}
    \item \textbf{RST encodes local coherence}: Discourse structure reflects how information flows coherently within a document, not ideological emphasis.
    
    \item \textbf{Noise in prominence estimation}: Tree depth and nuclearity are noisy proxies for rhetorical importance; many factors (sentence length, parsing errors) introduce variance.
    
    \item \textbf{Coverage confound}: Most variation in DFI vectors comes from coverage differences, which dominates any structural signal.
\end{enumerate}


\section{Limitations}

\begin{enumerate}
    \item \textbf{English only}: The pipeline operates only on English-language articles.
    
    \item \textbf{Dataset scope}: Results are specific to the AllSides dataset; generalization to other corpora is untested.
    
    \item \textbf{Triplet alignment errors}: The unsupervised triplet-finding procedure has an estimated 1\% false-triplet rate.
    
    \item \textbf{RST parser limitations}: We enforce a 5-second timeout that excludes some documents; parser errors may introduce noise.
    
    \item \textbf{Binary task}: We frame bias as left-vs-right comparison relative to center. This simplifies the three-way ideological space.
    
    \item \textbf{Cross-triplet alignment}: DFI dimensions are not semantically aligned across triplets; the model learns distributional patterns rather than fact-specific rules.
\end{enumerate}


\section{Conclusion}

We investigated whether RST structural prominence predicts media bias, finding that \textbf{it does not}. When fact coverage is controlled, structural features achieve only chance-level accuracy. Instead, the dominant signal is \textbf{fact omission}: which facts sources choose to cover or omit relative to a centrist baseline.

Our best model---Random Forest with bipartite EDU matching that explicitly encodes coverage asymmetry---achieves \textbf{89.77\% accuracy} in distinguishing left-vs-center from right-vs-center samples. Analysis reveals that right-leaning sources systematically omit more facts that centrist sources cover, particularly in early (salient) positions.

This work provides a novel, interpretable lens for media bias analysis. Rather than relying on opaque lexical features, we can characterize bias in terms of selective fact coverage---a finding with implications for media literacy and automated bias detection systems.

\section*{Code Availability}

Code and data are available at: \url{https://github.com/itsadityasai/rhetorical-framing-auditor}

\bibliographystyle{acl_natbib}
\bibliography{custom}

\end{document}
